\bookchapter{פרק 44: התנסות מ-0 ל-1 — תהליך מלא של פרויקט מקצה לקצה}

\chapterintro{התיאוריה נגמרה — עכשיו הזמן לעשות. הפרק הזה איננו מדריך “תקלידו אחריי” פשוט, אלא מתודולוגיה מעשית של “איך חושבים, איך שואלים, איך מאמתים”. נציג שלושה פרויקטים ברמות קושי שונות ונראה את תהליך החשיבה המלא מהדרישות ועד המוצר המוגמר. את התהליך הזה אפשר ליישם בכל כלי שתבחרו.}

\booksection{א. פרויקט 1: אתר כלים אישי מינימלי (רמת מתחילים)}

\booksubsection{חשיבה על דרישות}

מטרה: ליצור עמוד אינטרנט של אתר כלים אישי, שמציג כניסות ל-3-4 כלים שימושיים קטנים. זה פרויקט צד-לקוח בלבד, בלי שרת ובלי מסד נתונים. הדגש הוא לא על מורכבות הכלים עצמם, אלא על החוויה של התהליך המלא “מדרישות ועד מוצר”.

לפי חמשת יסודות הדרישה: Who — שימוש אישי. What — עמוד שמציג כניסות לכלים. Why — ריכוז הכלים הנפוצים בעמוד אחד לגישה נוחה. How — קובץ יחיד HTML+CSS+JS, בלי פריימוורק. Constraint — השלמה תוך 30 דקות, בלי פריסה.

\booksubsection{ניסוח ההוראה}

על בסיס עקרון Think Before Coding, ההוראה שלכם צריכה לכלול מספיק מידע כדי שה-AI יבין את המטרה, אבל אין צורך לפרט כל פרט מימוש:

\begin{lstlisting}
צור אתר כלים אישי לפי הדרישות הבאות:
1. HTML בקובץ יחיד (CSS ו-JS מוטמעים), בלי פריימוורקים
2. ערכת נושא כהה, סטייל מודרני ומינימליסטי
3. הצג 4 כרטיסי גישה לכלים:
   - מקצר טקסטים
   - סופר מילים
   - בוחר צבעים
   - מחולל קודי QR
4. לכל כרטיס אייקון, כותרת ותיאור קצר
5. לכרטיסים יש אנימציית ריחוף
6. פריסה רספונסיבית, שמישה בנייד
7. פס כותרת למעלה
תן ישירות את קובץ ה-HTML המלא.
\end{lstlisting}

\booksubsection{קבלה ואיטרציה}

אחרי שה-AI מספק את הקוד, עוברים קבלה לפי שלוש שכבות הבדיקה של האימות האיטרטיבי. שכבה ראשונה — אימות פונקציונלי: האם הקובץ נפתח בדפדפן? האם כל 4 הכרטיסים מוצגים? האם האנימציה בריחוף עובדת תקין? שכבה שנייה — אימות גבולות: האם הפריסה תקינה בגודל מסך של טלפון? האם הכרטיסים עוברים שורה כשהחלון צר מאוד? שכבה שלישית — אימות ארכיטקטוני: האם הקוד קריא? האם ה-CSS מרוכז בתג style אחד? האם יש מיותרים בולטים?

אם רוצים לממש גם את הפונקציונליות של כרטיס כלים כלשהו (ולא רק את הכניסה אליו), אפשר להוסיף בהדרגה — בכל פעם מוסיפים את הפונקציונליות של כלי אחד, מאמתים, ואז מוסיפים את הבא. זהו יישום מעשי של אסטרטגיית “צעדים קטנים ומהירים”.

\booksection{ב. פרויקט 2: מערכת ניהול הוצאות במחזור מלא (רמת מתקדמים)}

\booksubsection{חשיבה על דרישות ו-PRD}

מטרה: אפליקציית ניהול הוצאות במחזור מלא, עם צד-לקוח, צד-שרת ומסד נתונים. הפרויקט הזה מורכב בהרבה מהראשון, ודורש את שיטות הנדסת הדרישות מפרק 40.

קודם כותבים PRD. לפי התבנית מסעיף 3 של פרק 40, כותבים סקירת פרויקט, רשימת פיצ'רים, תוכנית טכנולוגית, מודל נתונים ותוכנית פיתוח. בבחירת הפתרון הטכנולוגי מתחשבים בעקרון Simplicity First — בצד הלקוח React אבל בלי Redux או ניהול מצב מורכב אחר; בצד השרת Node.js אבל בלי מיקרו-שירותים; והמסד יהיה SQLite ולא PostgreSQL (בשלב הפיתוח).

תוכנית הפיתוח מתפרקת ל-5 שלבים: שלב 1 — שלד API של השרת (שרת Express + SQLite + ממשקי CRUD); שלב 2 — עמודי בסיס בצד הלקוח (רשימת רשומות + טופס הוספה); שלב 3 — חיבור קצה-לקצה (קריאות API + תצוגת נתונים); שלב 4 — קטגוריות וסינון; שלב 5 — סיכום חודשי וגרפים פשוטים. בסוף כל שלב יש בידיכם “משהו שעובד”.

\booksubsection{הגדרת קובץ זיכרון הפרויקט}

לפני שמתחילים לקודד, יוצרים קודם את קובץ זיכרון הפרויקט. זו הפרקטיקה המרכזית של הנדסת ההקשר — לוודא שה-AI יודע את רקע הפרויקט בתחילת כל שיחה.

\begin{lstlisting}
# פרויקט: אפליקציית full-stack לניהול הוצאות
## סקירה
אפליקציה לניהול הוצאות אישיות עם תמיכה ברישום הכנסות והוצאות, קטגוריות וסיכומים חודשיים.
## סטאק טכנולוגי
- Frontend: React + TypeScript + Vite
- Backend: Node.js + Express + TypeScript
- מסד נתונים: SQLite (דרך better-sqlite3)
- טיפוסים משותפים לפרונט ולבק
## הסכמי API
- נתיב בסיס: /api
- פורמט תשובה: { success: boolean, data: any, error: string }
- סגנון RESTful
## מודל נתונים
Record { id, amount, type, category, note, date }
## כללי קוד
- מצב Strict ב-TypeScript
- הזחה של 2 רווחים
- פונקציות עם תיעוד JSDoc
- הערות והודעות commit בעברית
\end{lstlisting}

\booksubsection{מימוש שלב אחר שלב}

לפי תוכנית הפיתוח, כל שלב מתבצע בשיחה נפרדת משלו. לפניכם דוגמת הוראה לשלב הראשון:

\begin{lstlisting}
אני צריך ליצור שלד ל-API בבק.
דרישות:
1. שרת Express, פורט 3001
2. מסד נתונים SQLite, שם הטבלה records
3. שדות הטבלה records:
   id (TEXT PRIMARY KEY)
   amount (REAL)
   type (TEXT: 'income'|'expense')
   category (TEXT)
   note (TEXT)
   date (TEXT: YYYY-MM-DD)
4. נקודות קצה CRUD:
   GET /api/records - לקבל את כל הרשומות
   POST /api/records - להוסיף רשומה
   PUT /api/records/:id - לעדכן רשומה
   DELETE /api/records/:id - למחוק רשומה
5. מידלוור לטיפול בשגיאות
תן לי את מבנה הקוד המלא של הבק.
\end{lstlisting}

שימו לב למאפייני ההוראה הזאת: היא מתארת רק את השלב הנוכחי (שלד השרת) ולא מערבבת דרישות של צד הלקוח; היא מפרטת בבהירות את השדות והממשקים (נותנת ל-AI מספיק אילוצים), אבל את המימוש הספציפי מחליט ה-AI; והיא עוקבת אחרי עקרון Goal-Driven Execution — “מה לעשות” ולא “איך לעשות”.

אחרי שה-AI מספק את הקוד, מאמתים מיד: מריצים את השרת ובודקים עם curl או Postman שכל ממשק עובד כמו שצריך. אחרי אישור, פותחים שיחה חדשה לשלב 2 (עמודי בסיס בצד הלקוח). בכל שלב חוזרים על מעגל “הוראה-אימות-שיחה חדשה”.

\booksection{ג. פרויקט 3: אפליקציה מועצמת AI (רמה גבוהה)}

\booksubsection{גיבוש רעיון הפרויקט}

מטרה: אפליקציה שמתחברת ל-API של מודל גדול — למשל “עוזר למידה AI”: המשתמש מזין חומר לימוד, וה-AI מייצר סיכום, מקשר בין נושאים וממליץ על מקורות למידה. הייחוד של הפרויקט הזה: האפליקציה שלכם עצמה קוראת ל-AI, ואתם צריכים להתמודד עם בעיות חדשות כמו קריאות API, הנדסת Prompt וטיפול בשגיאות.

הלקח החשוב ביותר שהפרויקט הזה מלמד: כש“משתמשים ב-AI כדי לכתוב קוד שקורא ל-AI”, צריך לשים לב במיוחד להנדסת ה-Prompt. גם ה-Prompt שהאפליקציה שלכם משתמשת בו בקריאה ל-API של המודל הגדול צריך לעמוד בעקרונות שדנו בהם במדריך הזה — בהירות, תמציתיות ומיקוד במטרה.

\booksubsection{נקודות טכניות מרכזיות}

הפרויקט הזה כולל כמה נקודות טכניות שלא פגשתם בשני הפרויקטים הקודמים:

\begin{itemize}
  \item ניהול מפתחות API: אסור לקודד את המפתח ישירות בקוד — מנהלים אותו במשתני סביבה;
  \item טיפול בתשובות זורמות: התשובות של המודל הגדול בדרך כלל זורמות, וצריך להציג אותן מילה אחר מילה;
  \item שליטה בצריכת Token: המודל הגדול מחויב לפי Token, וצריך לשלוט בצריכה ברמת האפליקציה;
  \item ניהול תבניות Prompt: לכל פונקציונליות תבנית Prompt משלה, וצריך לנהל אותן בצורה מובנית.
\end{itemize}

את כל הנקודות הטכניות האלה אפשר לבקש מה-AI לממש בשבילכם, אבל צריך להבין את ה“למה” של כל אחת מהן — למה אסור לקודד מפתחות בקוד? למה צריך תשובות זורמות? למה צריך לשלוט בצריכת Token? רק כשמבינים את ה“למה” אפשר לשפוט אם הפתרון שה-AI הציע הגיוני.

\booksection{ד. צ'ק-ליסט אוניברסלי לתהליך המעשי}

בין אם הפרויקט גדול ובין אם קטן, לפניכם צ'ק-ליסט אוניברסלי לתהליך תכנות AI. בכל פעם שמתחילים פרויקט חדש, עברו על הרשימה וודאו שלא פספסתם שלב מרכזי.

\begin{center}
\begin{tabularx}{\linewidth}{|X|X|X|}
\hline
\textbf{שלב} & \textbf{מה עושים} & \textbf{העיקרון המתאים} \tabularnewline
\hline
1. חשיבה על דרישות & עוברים על חמשת יסודות הדרישה ומבהירים מה רוצים לבנות & Think Before Coding \tabularnewline
2. פירוק הדרישות & מפרקים דרישה גדולה לשלבים קטנים שאפשר לאמת בנפרד & Simplicity First \tabularnewline
3. כתיבת PRD & מארגנים למסמך מובנה שישמש לעיון & Goal-Driven \tabularnewline
4. הגדרת קובץ זיכרון & יוצרים קובץ CLAUDE.md או קובץ דומה & Harness Engineering \tabularnewline
5. מימוש שלב אחר שלב & שיחה אחת לכל שלב: הוראה-אימות-המשך & Surgical Changes \tabularnewline
6. בדיקות ואימות & מיד אחרי כל שלב בודקים את תקינות הפונקציונליות & Iterate \& Verify \tabularnewline
7. סקירת קוד & בודקים הגיון ארכיטקטוני, אבטחה ותחזוקתיות & Iterate \& Verify \tabularnewline
8. איטרציה ושיפור & משפרים באיטרציות לפי תוצאות הבדיקות והסקירה & Iterate \& Verify \tabularnewline
\hline
\end{tabularx}
\end{center}

\begin{center}\small\color{gray}צ'ק-ליסט אוניברסלי לתהליך תכנות AI\end{center}

\begin{tcbpractice}{הידיים על המקלדת: השלימו את “פרויקט 1” השבוע}
התייחסו ל“אתר הכלים האישי” מסעיף 1 של פרק 44 כאל מטלת תכנות ה-AI הראשונה שלכם: הקדישו 5 דקות לחמשת היסודות, שלחו את ההוראה ל-AI (כל כלי יספיק), ואז קבלו לפי שלוש שכבות האימות. סיימו אותו תוך 30 דקות — והינה נולדה יצירת תכנות ה-AI הראשונה שלכם. אחרי הגמר אל תשכחו לשמור אותו היטב — זו המזכרת שלכם על המעבר מ“קורא” ל“עושה”.

\end{tcbpractice}

\begin{tcbsummary}{סיכום הפרק}
\begin{itemize}
  \item פרויקט 1 (מתחילים): אתר כלים בצד הלקוח בלבד — חוויה של התהליך המלא, נגמר תוך 30 דקות
  \item פרויקט 2 (מתקדמים): מערכת ניהול הוצאות במחזור מלא — PRD + קובץ זיכרון + מימוש ב-5 שלבים, ובכל שלב “משהו שעובד”
  \item פרויקט 3 (רמה גבוהה): עוזר למידה AI — לכתוב עם AI קוד שקורא ל-AI, תוך תשומת לב להנדסת Prompt
  \item תהליך אוניברסלי בשמונה שלבים: חשיבה על דרישות פירוק PRD קובץ זיכרון מימוש הדרגתי בדיקות סקירה איטרציה
  \item בכל שלב עוקבים אחרי מעגל “הוראה-אימות-שיחה חדשה”
\end{itemize}

\end{tcbsummary}
